\section{Localizing the entropy inequality to a prescribed geodesic}
\label{sec:every-geodesic}

Proposition~\ref{weight:variable-cd} provides an entropy inequality
for optimal transports between measures. To recover a statement on
one prescribed geodesic, we concentrate the endpoint measures in
shrinking balls. Nonbranching makes every strict interior subsegment
unique, so all the resulting transport paths converge to that
subsegment. At the same time, we increase \(\lambda\) fast enough
that endpoint entropy contributes no error in the limit.

\begin{proof}[Proof of Theorem~\ref{weight:every-geodesic}]
The case \(\ell=0\) is immediate. A finite-dimensional
\(\RCD\) space is nonbranching \cite[Theorem~1.3]{Deng2020}.
Every strict subsegment of \(\sigma\) is therefore the unique
minimizing geodesic between its endpoints. Indeed, replacing
such a subsegment by a different minimizer gives a minimizing
curve between \(\sigma(0)\) and \(\sigma(1)\), with the same
length \(\ell\), agreeing with \(\sigma\) on a nontrivial
initial interval. This contradicts nonbranching.

Fix one strict subsegment, and denote its constant-speed
reparametrization by \(\tau\colon[0,1]\to M\).
For \(\varepsilon_j\downarrow0\), let
\[
\mu_i^j=\frac{m|_{B_{\varepsilon_j}(\tau(i))}}
                  {m(B_{\varepsilon_j}(\tau(i)))},\qquad
E_i^j=\Ent_m(\mu_i^j)
     =-\log m(B_{\varepsilon_j}(\tau(i))),\quad i=0,1.
\]
These quantities are finite by full support and local finiteness.
Choose
\[
\lambda_j\ge j\bigl(1+|E_0^j|+|E_1^j|\bigr).
\]
The two endpoint measures have finite entropy relative to
\(m_{\lambda_j}\), so Proposition~\ref{weight:variable-cd}
provides a dynamical optimal plan \(\pi_j\) satisfying the
variable-curvature entropy inequality. Write
\(\mu_t^j=(e_t)_\#\pi_j\).

All geodesics in these plans lie in one compact set \(Q\) and
have lengths bounded by one constant \(L\). Properness and
Arzel\`a--Ascoli give compactness of this family of paths.
Every limit of a path with these shrinking endpoint constraints
must be \(\tau\), by the uniqueness just proved. Hence
\(\pi_j\) converges weakly to \(\delta_\tau\), even though
the measures \(m_{\lambda_j}\) vary with \(j\).

Set \(g(s,t)=\min\{s,t\}-st\), the nonnegative Dirichlet
Green kernel on \([0,1]\), and define
\begin{align*}
A_t^j&=\int F\dd\mu_t^j,& E_t^j&=\Ent_m(\mu_t^j),\\
C_t^j&=\int_{\Geo(M)}\int_0^1g(s,t)|\dot\gamma|^2
                      \dd s\dd\pi_j(\gamma),\\
D_t^j&=\int_{\Geo(M)}\int_0^1g(s,t)G(\gamma_s)|\dot\gamma|^2
                      \dd s\dd\pi_j(\gamma).
\end{align*}
We use the large-parameter entropy comparison appearing in
\cite[Theorem~3.13, Equation~(3.6)]{Han2018Monotonicity}, now with the variable
bound and the concentrating endpoint measures chosen above.
Using \(\Ent_{m_{\lambda_j}}(\mu)=\Ent_m(\mu)
-\lambda_j\int F\dd\mu\), the defining
\(\CD(K-\lambda_jG,\infty)\) inequality becomes
\begin{equation}
\label{weight:entropy-rearranged}
A_t^j\ge(1-t)A_0^j+tA_1^j-D_t^j
 +\frac{E_t^j-(1-t)E_0^j-tE_1^j+KC_t^j}{\lambda_j}.
\end{equation}
This uses the Green-kernel normalization of
\cite[Definition~1.2]{BHS2021}. The intermediate entropies
are finite, and Jensen's inequality on \(Q\) gives
\(E_t^j\ge-\log m(Q)\). Also
\(0\le C_t^j\le L^2/8\).
Thus the last term in \eqref{weight:entropy-rearranged} has
lower limit at least zero, uniformly in \(t\).
The continuity of \(F,G\) on \(Q\), compactness of the paths,
and \(\pi_j\to\delta_\tau\) imply
\begin{equation}
\label{weight:green-concavity}
F(\tau_t)\ge(1-t)F(\tau_0)+tF(\tau_1)
-|\dot\tau|^2\int_0^1g(s,t)G(\tau_s)\dd s.
\end{equation}

The argument applies to every strict subsegment of \(\sigma\).
On any compact interval inside \((0,1)\), choose a twice
continuously differentiable function \(a\) with
\(a''(t)=\ell^2G(\sigma_t)\).
The Green formula and \eqref{weight:green-concavity}, applied
on all subintervals, say exactly that \(F\circ\sigma-a\)
lies above each of its chords. It is concave, so its second
distributional derivative is nonpositive. This proves
\eqref{weight:curve-hessian}. Continuity supplies the values
at the ends of any closed subinterval.
\end{proof}

